%-----------------------------------------------------------------------------------------------------------------------------------------------%
%	The MIT License (MIT)
%
%	Copyright (c) 2021 Jitin Nair
%
%	Permission is hereby granted, free of charge, to any person obtaining a copy
%	of this software and associated documentation files (the "Software"), to deal
%	in the Software without restriction, including without limitation the rights
%	to use, copy, modify, merge, publish, distribute, sublicense, and/or sell
%	copies of the Software, and to permit persons to whom the Software is
%	furnished to do so, subject to the following conditions:
%	
%	THE SOFTWARE IS PROVIDED "AS IS", WITHOUT WARRANTY OF ANY KIND, EXPRESS OR
%	IMPLIED, INCLUDING BUT NOT LIMITED TO THE WARRANTIES OF MERCHANTABILITY,
%	FITNESS FOR A PARTICULAR PURPOSE AND NONINFRINGEMENT. IN NO EVENT SHALL THE
%	AUTHORS OR COPYRIGHT HOLDERS BE LIABLE FOR ANY CLAIM, DAMAGES OR OTHER
%	LIABILITY, WHETHER IN AN ACTION OF CONTRACT, TORT OR OTHERWISE, ARISING FROM,
%	OUT OF OR IN CONNECTION WITH THE SOFTWARE OR THE USE OR OTHER DEALINGS IN
%	THE SOFTWARE.
%	
%
%-----------------------------------------------------------------------------------------------------------------------------------------------%

%----------------------------------------------------------------------------------------
%	DOCUMENT DEFINITION
%----------------------------------------------------------------------------------------

% article class because we want to fully customize the page and not use a cv template
\documentclass[a4paper,10pt]{article}

%----------------------------------------------------------------------------------------
%	PACKAGES
%----------------------------------------------------------------------------------------
\usepackage{url}
\usepackage{parskip} 	

%other packages for formatting
\RequirePackage{color}
\RequirePackage{graphicx}
\usepackage[usenames,dvipsnames]{xcolor}
\usepackage[scale=0.9]{geometry}

%tabularx environment
\usepackage{tabularx}

%for lists within experience section
\usepackage{enumitem}

% centered version of 'X' col. type
\newcolumntype{C}{>{\centering\arraybackslash}X} 

%to prevent spillover of tabular into next pages
\usepackage{supertabular}
\usepackage{tabularx}
\newlength{\fullcollw}
\setlength{\fullcollw}{0.47\textwidth}

%custom \section
\usepackage{titlesec}				
\usepackage{multicol}
\usepackage{multirow}

%CV Sections inspired by: 
%http://stefano.italians.nl/archives/26
\titleformat{\section}{\large\scshape\raggedright}{}{0em}{}[\titlerule]
\titlespacing{\section}{0pt}{6pt}{6pt}

%for publications
\usepackage[style=authoryear,sorting=ynt, maxbibnames=2]{biblatex}

%Setup hyperref package, and colours for links
\usepackage[unicode, draft=false]{hyperref}
\definecolor{linkcolour}{rgb}{0,0.2,0.6}
\hypersetup{colorlinks,breaklinks,urlcolor=linkcolour,linkcolor=linkcolour}
\addbibresource{citations.bib}
\setlength\bibitemsep{1em}

%for social icons
\usepackage{fontawesome5}
\renewcommand{\familydefault}{\sfdefault}

%debug page outer frames
%\usepackage{showframe}


% job listing environments
\newenvironment{jobshort}[2]
    {
    \begin{tabularx}{\linewidth}{@{}l X r@{}}
    \textbf{#1} & \hfill &  #2 \\
    \end{tabularx}
    \vspace{-4pt}
    }
    {
    \vspace{3pt}
    }

\newenvironment{joblong}[2]
    {
    \begin{tabularx}{\linewidth}{@{}l X r@{}}
    \textbf{#1} & \hfill &  #2 \\[3.75pt]
    \end{tabularx}
    \begin{minipage}[t]{\linewidth}
    \begin{itemize}[nosep,after=\strut, leftmargin=1em, itemsep=3pt,label=--]
    }
    {
    \end{itemize}
    \end{minipage}    
    }


%----------------------------------------------------------------------------------------
%	BEGIN DOCUMENT
%----------------------------------------------------------------------------------------
\begin{document}

% non-numbered pages
\pagestyle{empty} 

%----------------------------------------------------------------------------------------
%	TITLE
%----------------------------------------------------------------------------------------
\begin{center}
\LARGE{June (Junhyeok) Lee}\\[7.5pt]
\normalsize{
\href{https://github.com/born-june04}{\raisebox{-0.05\height}\faGithub\ born-june04} \ $|$ \ 
\href{https://linkedin.com/in/junh-lee/}{\raisebox{-0.05\height}\faLinkedin\ junh-lee} \ $|$ \ 
\href{https://born-june04.github.io/}{\raisebox{-0.05\height}\faGlobe \ Portfolio} \ $|$ \ 
\href{mailto:june0604@uw.edu}{\raisebox{-0.05\height}\faEnvelope \ june0604@uw.edu} \ $|$ \ 
\href{tel:+12065127412}{\raisebox{-0.05\height}\faMobile \ +1.206.512.7412}
}\\[5pt]
\end{center}

%Interests/ Keywords/ Summary
\section{Summary}
I develop clinically aligned representation learning systems for biomedical AI. I move beyond black-box optimization by embedding physiological and structural priors into models to improve robustness under real-world distribution shift.


%----------------------------------------------------------------------------------------
%	EDUCATION
%----------------------------------------------------------------------------------------
\section{Education}
\begin{tabularx}{\linewidth}{@{}l X@{}}
2025 -- Present & \textbf{M.S. in Bioengineering}, University of Washington \hfill GPA: 3.9/4.0 \\
& Focus: Multi-scale biomechanics and physics-informed deep learning \\
& \textit{Expected Graduation: June 2026} \\[4pt]

2018 -- 2024 & \textbf{B.S. in Biomedical Engineering}, Hankuk University of Foreign Studies \hfill GPA: 3.9/4.0 \\
& Thesis: \textit{Deep Clustering of Single-cell RNA-seq with Combined Features} \\
\end{tabularx}


%----------------------------------------------------------------------------------------
%	WORK EXPERIENCE
%----------------------------------------------------------------------------------------
%Experience
\section{Work Experience}

\begin{joblong}{Research Scientist | Monitor Corporation}{Jul 2023 -- Jun 2025}
\item \textbf{Probe-Conditional Ultrasound Denoising (Production DDPM)}:
\begin{itemize}[nosep,after=\strut, leftmargin=2em, itemsep=1pt,label=$\circ$]
\item Reduced false positives by \textbf{89\% across 50K+ clinical scans} by deploying a probe-conditional DDPM with contrastive learning to suppress spectral artifacts across three GE ultrasound probe variants
\item Eliminated per-probe retraining by unifying domain adaptation into a single conditional architecture, explicitly conditioning on probe-specific prior distributions to address hardware-induced domain shift
\end{itemize}
\item \textbf{Semi-Supervised 3D Segmentation Pipeline}:
\begin{itemize}[nosep,after=\strut, leftmargin=2em, itemsep=1pt,label=$\circ$]
\item Achieved \textbf{0.72 Dice} with sub-second inference for tumor segmentation by building a Fast Fourier Convolution (FFC) 3D U-Net with multi-scale frequency filtering
\item Reduced annotation cost by \textbf{21,600 hours (\$540K)} by designing a MedSAM-inspired semi-supervised framework with reliability scoring, scaling \textbf{300 to 21,900 labeled scans} (1:73 ratio) while maintaining fully supervised parity
\end{itemize}
\item \textbf{Model Compression for Edge Deployment}:
\begin{itemize}[nosep,after=\strut, leftmargin=2em, itemsep=1pt,label=$\circ$]
\item Achieved \textbf{2.5$\times$ faster inference} preserving \textbf{98\% accuracy} by applying bfloat16 quantization and structured pruning to compress 3D segmentation models by 50\%, enabling real-time clinical deployment
\end{itemize}
\end{joblong}

%----------------------------------------------------------------------------------------
%	RESEARCH EXPERIENCE
%----------------------------------------------------------------------------------------
\vspace{-7pt}
\section{Research Experience}

\begin{joblong}{Vector Decomposition from Frozen Segmentation Encoders | Asan Hospital }{Jan 2025 -- Present}
\item Discovered latent biomarker axes in frozen segmentation encoders recoverable with \textbf{zero trainable parameters}, achieving directional stability (5-fold cos$\theta$=0.90/0.96) via centroid-difference projections
\item Demonstrated zero-shot generalization to held-out patients (\textbf{CVS AUC 0.803, PRL AUC 0.716}) without any task-specific training
\item Validated physical grounding with imaging gradients (SWI gradient $p<0.001$; PRL rim $\rho=-0.49$, $p<0.001$); submitted to \textit{MICCAI 2026} (2nd author)
\end{joblong}

\vspace{-5pt}

\begin{joblong}{WiseFind: Structure-as-Context Multimodal RAG | Harborview Medical Center (HMC)}{Sep 2025 -- Present}
\item Identified intra-domain visual homogeneity as a fundamental retrieval barrier (silhouette $<$0.06 post-fine-tuning), motivating text- and structure-based retrieval over visual embeddings
\item Achieved \textbf{P@1 0.848, MRR 0.872, R@5 0.701}\textbf{---3.3$\times$} above best visual baseline on 200 expert neurosurgical queries by proposing Shotgun Retrieval via THC Tree with three complementary channels (T2I, CI, XRef)
\item Raised hard-query P@1 by \textbf{+76\%} for figures $>$10 pages from text anchor by fine-tuning ColBERT (PubMedBERT backbone, 30K steps) on 5,544 domain-specific pairs with cross-reference following
\item Submitted to \textit{MICCAI 2026} (1st author)
\end{joblong}

\vspace{-5pt}

\begin{joblong}{WiseMind: RAG-Powered Neurosurgical AI Assistant | HMC}{Feb 2026 -- Present}
\item Reduced hallucination in medical LLM responses by designing a \textbf{multi-step chain-of-thought pipeline} (Draft $\rightarrow$ Review $\rightarrow$ Refine) with retrieval-grounded inline \textbf{[N] citations} linking each claim to exact source passages for clinician verification
\item Improved clinical recall by engineering \textbf{ontology-driven query expansion} (1,730 medical terms) and domain-scoped retrieval with automatic fallback, covering \textbf{27,845 indexed passages} across textbook and guideline corpora
\item Achieved source-faithful multi-document synthesis by prompting \textbf{MedGemma 27B} to present conflicting evidence side-by-side rather than silently resolving disagreements between clinical guidelines
\item Building an \textbf{RLHF data pipeline} from clinician thumbs-up/down feedback and RAG-provenance-annotated conversations stored in MongoDB, targeting reward model training for domain-aligned LLM fine-tuning; deployed at HMC with 7 release iterations (v1.0--v1.4)
\end{joblong}

\vspace{-5pt}

\begin{joblong}{WiseSpine: Patient-Specific Vertebral Fracture Simulation | HMC}{Sep 2025 -- Present}
\item Reproduced \textbf{full A0--A4 AO classification} from first principles by building a patient-specific voxel finite-element engine (234K hexahedral elements, $Ku{=}F$) with HU-derived anisotropic material properties directly from clinical CT
\item Achieved \textbf{clinically consistent force--fracture thresholds} (A0 at 1\,kN $\rightarrow$ A4 at 8\,kN; BMD\,0.3 $\rightarrow$ 77.5\% yield) by implementing transversely isotropic elasticity, incremental loading, and progressive continuum damage mechanics with cortical shell blending
\item Enabled \textbf{per-scenario simulation in $<$20\,min} on a single GPU by accelerating sparse assembly and CuPy-based conjugate gradient solve over 784K DOF, making patient-specific fracture risk assessment clinically feasible
\item Formulated segmentation robustness as a \textbf{physics-constrained counterfactual augmentation} framework with causal DAG analysis (force, BMD, flexion $\rightarrow$ AO type), generating balanced A0--A4 training data from normal CT scans; preparing \textit{NeurIPS 2026}
\end{joblong}

\vspace{-5pt}

\begin{joblong}{Culture-Positive TB Detection and Physiologically Grounded AI Modeling | Ubicomp Lab}{Sep 2025 -- Present}
\item Reduced TB diagnostic latency from weeks to real-time, achieving \textbf{AUC 0.9519} and \textbf{92.7\% sensitivity} by developing an acoustic detection system on 7,031 cough segments
\item Identified infection-linked spectral signatures by extracting \textbf{57 physiological biomarkers}, revealing demographic and bacterial-load heterogeneity (female-specific slope $d=0.593$; age-dependent reversals)
\item Demonstrated interpretability via saliency concentration (25.9\% in 500--3000 Hz) and biomarker correlation ($|r|$ up to 0.42); in preparation for \textit{Nature Communications}
\end{joblong}

\vspace{-5pt}

\begin{joblong}{Tissue-Aware Drug Response Prediction | Bioinformatics Lab}{Jan 2022 -- Dec 2025}
\item Achieved \textbf{$R^2 = 0.857 \pm 0.012$} with \textbf{61.6\% MSE reduction} on 257,955 drug--cell line pairs by modeling tissue identity as a structural prior using 54,837 GTEx samples
\item Validated biological grounding of gene-importance estimation (DrugBank/KEGG $p < 10^{-8}$; $r=0.40$) by designing a cross-attention multimodal architecture
\item Manuscript under review at \textit{Bioinformatics}
\end{joblong}

\vspace{-5pt}

\begin{joblong}{MRI Bias Field Correction with Generative AI | MedicalVision Lab}{Aug 2022 -- Apr 2023}
\item Achieved \textbf{9.75$\times$ faster inference} than standard DDPM while preserving diagnostic fidelity by developing a diffusion-based MRI bias field correction with Hampel mixture noise modeling; published at \textbf{\textit{MIDL 2023}}
\end{joblong}


\section{Hackathon}

\begin{joblong}{Google MedGemma 2026 Healthcare Hackathon \textnormal{(\href{https://github.com/born-june04/medgemma_2026challenge}{GitHub}, \href{https://www.youtube.com/watch?v=j5Z_kRi3wDU}{Demo})}}{Jan 2026 -- Present}
\item Achieved end-to-end multi-modal clinical diagnosis in $\sim$5s on consumer hardware by architecting a \textbf{three-agent pipeline} (Audio/HeAR, CXR/MedSigLIP, CT/MedSigLIP) with confidence-weighted late fusion mirroring the multi-specialist case conference model
\item Produced interpretable \textit{biological stories} linking findings to diagnoses by designing a \textbf{3-level hierarchical reasoning framework} (pathophysiology $\rightarrow$ pattern recognition $\rightarrow$ disease-specific biomarkers) with explicit quantitative thresholds
\item Formalized \textbf{disagreement as a clinical signal}: divergent agents raise a \textit{NEEDS CONFIRMATION} flag with modality-specific weights (CT $\times$1.2, CXR $\times$1.0, Audio $\times$0.8), surfacing co-morbidity hypotheses rather than silently overriding minority agents
\item Designed for \textbf{modular deployment in resource-limited settings} --- audio-only smartphone screening (HeAR) gracefully scaling to CXR then CT; FHIR-compatible JSON outputs; submitted to \textit{MedGemma Impact Challenge 2026} (Kaggle)
\end{joblong}


\begin{joblong}{Seoul National University Hospital (SNUH) Skin Burn Classification Hackathon --- 3rd Place}{2022}
\item Improved burn depth classification without increasing model complexity by introducing a physiologically motivated attenuation-inspired RGB reweighting scheme (0.21, 0.72, 0.07), emphasizing hemoglobin-sensitive green-channel information
\item Awarded \textbf{3rd Place} among national-level medical AI teams
\end{joblong}

%----------------------------------------------------------------------------------------
%	PUBLICATIONS
%----------------------------------------------------------------------------------------
\section{Publications}

\textbf{Lee, Junhyeok}, Alex Ching*, Cynthia Dong*, David Horne, Thomas Hawn, Shwetak Patel. (2026). Rapid Identification of Culture-Positive Tuberculosis from Cough Sounds Using Deep Learning. \textit{Nature Communications} (Manuscript in preparation).

\textbf{Lee, Junhyeok}, Aruna Souri*, Mary Nguyen*, Benjamin Han*, Christopher W. Lewis, Diana Wiseman. (2026). WiseFind: Why Structure Matters for Traceable Retrieval in Neurosurgery. \textit{Medical Image Computing and Computer Assisted Intervention (MICCAI)} (under review).

Jeong, Eunseon, \textbf{Lee, Junhyeok}, Yangshin Choi, Yoonho Nam. (2026). Extracting CVS and PRL Biomarker Axes from Frozen Lesion Segmentation Encoders. \textit{Medical Image Computing and Computer Assisted Intervention (MICCAI)} (under review).

\textbf{Lee, Junhyeok}, Jiwon Seo, Younhee Ko. (2026). Interpretable Tissue-Aware Drug Response Prediction Through Biologically Grounded Multimodal Deep Learning. \textit{Bioinformatics} (under review).

\textbf{Lee, Junhyeok*}, J. Kang*, Y. Nam, and T. Lee. (Apr. 2023). Bias Field Correction in MRI with Hampel Noise Denoising Diffusion Probabilistic Model. \textit{Proceedings of Medical Imaging with Deep Learning (MIDL)}.


%----------------------------------------------------------------------------------------
%	TEACHING
%----------------------------------------------------------------------------------------
\section{Teaching Experience}

\begin{jobshort}{Teaching Assistant | BIOE 437/537 Computational Systems Biology}{Fall 2025}
\textit{University of Washington, Prof. Herbert M. Sauro \& Prof. Joseph L. Hellerstein}
\end{jobshort}
\begin{itemize}[nosep,after=\strut, leftmargin=1em, itemsep=2pt,label=--]
\item Assisting enzyme kinetic simulation and optimization coursework; leading office hours and grading
\item Contributing to novel research integrating simulation-generated data for AI training
\end{itemize}
\vspace{2.0pt}


%----------------------------------------------------------------------------------------
%	SKILLS
%----------------------------------------------------------------------------------------
\section{Technical Skills}
\begin{tabularx}{\linewidth}{@{}l X@{}}
Programming Languages &  \normalsize{Python, R, Matlab, SQL, C/C++, Flask, SvelteKit, Neo4j, MongoDB, ChromaDB}\\
ML/AI Frameworks &  \normalsize{PyTorch, Transformers, DDPM, GenAI, AWS SageMaker, LLM, RAG, ColBERT, FAISS, Ollama}\\
Systems/Infrastructure &  \normalsize{PyCuda, TensorRT, Distributed Training, Docker, SLURM, HPC Cluster, Apptainer, nginx, Gunicorn}\\
Project Management &  \normalsize{Atlassian Jira, Confluence, Github}\\
Research Methodologies &  \normalsize{Physical Simulation, Bayesian Optimization, Statistical Validation, Clinical Trials}\\
\end{tabularx}
\vspace{2.0pt}

%----------------------------------------------------------------------------------------
%	AWARDS
%----------------------------------------------------------------------------------------
\section{Awards \& Honors}
\begin{tabularx}{\linewidth}{@{}l X r@{}}
2022 & \textbf{Director's Medal}, Korean Information Society Development Institute & ACK \\
2022 & \textbf{3rd Place}, Skin Burn Classification Hackathon & Seoul National University Hospital \\
2021 & \textbf{Army Commendation Medal (ARCOM)} & U.S. Army \\
\end{tabularx}


\vfill
\center{\footnotesize Last updated: \today}

\end{document}